\documentclass[11pt]{article}
\usepackage[a4paper,margin=25mm]{geometry}
\usepackage[T1]{fontenc}
\usepackage[utf8]{inputenc}
\usepackage{lmodern}
\usepackage{microtype}
\usepackage{booktabs}
\usepackage{tabularx}
\usepackage{enumitem}
\usepackage[round,authoryear]{natbib}
\usepackage[hidelinks]{hyperref}
\setlist{nosep}
\newcommand{\tbd}[1]{\textbf{[#1]}}

\title{TRACE: How studies of ``AI fatigue'' tie fatigue to AI, measure it and report it.\\
A registered methodological review}
\author{Mikoto Miura (Independent Researcher) \and \tbd{Author 2: name, affiliation}}
\date{Draft for discussion between the authors, 17 October 2026. Not for circulation.}

\begin{document}
\maketitle

\begin{quote}\small
\textbf{Status of this draft.} The protocol is not yet registered, and no record of the main set has been coded. The
Introduction, Background and Method follow registration draft 0.10 and the manuscript outline (version 3). The Results
and Discussion contain only the headings, the table shells and what each part will report; they contain no result and
no expectation about results. Items in bold brackets are to be agreed between the authors.
\end{quote}

\begin{abstract}
\noindent\textbf{Background.} Work on AI-related fatigue is growing; reviews call for validated measures and place it
within digital burnout.
\textbf{Question.} How do such studies tie fatigue to AI, and do their titles and conclusions stay within that tie,
their design and their instruments' names?
\textbf{Methods.} TRACE, a registered methodological review, coded studies from one free-source literature database,
screened by authors who knew the rules and closed on 31 December 2026 (not a census). \tbd{Coding as it happened: two
coders sharing the view of an earlier narrative review / the first author and a third person / the first author twice}
applied seven codes.
\textbf{Results.} \tbd{Written from the tables after coding, by the registered abstract rule.}
\textbf{Availability.} Registration, outline, code table with per-instrument and per-study tables, codebook and scripts
are open.
\end{abstract}

\section{Introduction}

This review covers studies that measure fatigue in relation to the use of AI systems, whether or not they call it
``AI fatigue''. Under that term the literature places constructs that range from fatigue attributed to AI in the items
of a scale to general fatigue measured among people who use AI. Work on the topic is growing
\citep{ragolane2025}. One review calls for standardized, validated measures of AI fatigue \citep{sosinski2026};
another places AI fatigue within digital burnout, with which it shares symptoms \citep{ragolane2025}.

A conclusion about ``AI fatigue'' depends on three things: whether the measure ties fatigue to AI, what the fatigue
items ask about, and whether what the report states is internally consistent and matches its conclusion. Item wording
shows what an instrument asks about. It does not show what the scores contain, which is a question of structural and
construct validity and needs correlational or factor-analytic data. This review counts motivational and behavioural
depletion as fatigue, so items about reduced interest, avoidance or irritation are described rather than treated as
outside the construct. Overlapping samples change how much evidence a set of reports carries.

Measurement research in general shows that a measure's name does not tell what it measures, and that items are
dropped, added and altered while the instrument keeps the same name; a methodological review is one way to make such
flexibility visible \citep{elson2023}. Studies of measurement practice in other fields count what reports disclose
\citep{perrig2024}, and reviews of research on trust in AI have counted how often studies built their own
questionnaire items rather than using existing questionnaires \citep{bach2022,benk2024}. For AI-related instruments,
systematic reviews have assessed measurement properties with COSMIN \citep[for example,][on AI literacy
scales]{lintner2024}, and a registered review of AI anxiety and attitude scales has no results yet
\citep{jonaid2026}.

We did not find a review of AI-fatigue studies that tabulates, study by study, how fatigue is tied to AI, what the
fatigue items ask about, whether the reports are internally consistent, and whether titles and conclusions stay within
the level of attribution, the design and the name of the instrument. For internal consistency, 16 searches of free
sources (Europe PMC, OpenAlex and a web search engine, 17 October 2026) found no such review; for the other points, the
claim rests only on the reviews we had read.

Our earlier critical narrative review \tbd{citation and authorship to be agreed} described measurement problems in 11
reports drawn from the same exploratory bibliography, without registration. We cite it as prior work and do not
restate its observations as findings. What this review adds is limited to five points: rules fixed and registered
before coding; primary counts from records that were not used to build the rules; item-level coding of instrument
content with targets; codes that the narrative review does not apply study by study (the AI-side variable, the title
checks and the three conditions on the main conclusion); and recalculation and consistency rules applied to every
record. The observations of the narrative review are set beside our statements, not tested against them.

The review asks one question: when a study reports fatigue in relation to AI, how is the fatigue tied to AI, and do
its title and conclusion stay within that tie, the study's design and the name of its instrument? The registered
coding scheme, TRACE, answers it with checks of attribution, names and design, not with validity evidence. The
primary research questions are:
\begin{itemize}
  \item RQ1 (attribution): at which level do studies tie their main fatigue instrument to AI?
  \item RQ2 (instruments and items): how do fatigue instruments break down by origin and modification, reported
  adaptation procedures, reported measurement properties and item content, and what instruments are used for the
  AI-side variable?
  \item RQ3 (reporting): how complete and how internally consistent is reporting?
  \item RQ4 (sample overlap): how many records are linked as the same sample or show unconfirmed signs of overlap, and
  how many studies remain when linked records are counted once?
  \item RQ5 (title versus content): do titles and abstracts match the name of the main instrument, the design and the
  level of attribution?
\end{itemize}
The secondary questions are how these answers are distributed within each of the source database's two eligibility
groups (SQ1) and, for each of three conditions, whether a study's main conclusion meets it (SQ2).

The contributions are: (1) TRACE, a registered coding scheme that links, in each study, how fatigue is tied to AI,
what the fatigue items ask about, and what the report states and concludes; (2) a description of how a set of studies
not used to build the rules attributes fatigue to AI, measures it and reports it; (3) the item-level content of the
fatigue instruments; and (4) recommendations for authors and reviewers that separate what existing standards already
ask from steps that the standards cited here do not name. The released code table, codebook and scripts let the tables
be recomputed and the scheme be reused; they are not claimed as findings.

\section{Background}

\subsection{Fatigue, exhaustion and burnout}
The eligibility of this review defines fatigue to include exhaustion, burnout and motivational and behavioural
depletion \tbd{the definition of the source database, quoted from the registration}. The definition shapes two
outcomes: the share of items about reduced interest is partly set at entry, and the value ``no attribution to AI'' is
truncated because studies without any relation to AI are excluded.

\subsection{Technostress, adapted scales and scales built for AI}
\tbd{Technostress and its exhaustion component; general fatigue scales reworded for AI; scales built for AI. Sources to
be read in free full text and registered before they are cited.}

\subsection{Measurement flexibility and names}
A measure's name is often read as a statement about its content, and instruments are changed while their names stay the
same \citep{elson2023}. Four codes of this review apply these known problems to the AI-fatigue literature: the
modification of instruments (C2a), what reports disclose about them (C2d), and two checks that compare names (C6a for
the title and the third condition of C7 for the main conclusion). The name checks compare names; they do not compare
the content of instruments.

\subsection{Content validity and reviews of measurement properties}
COSMIN evaluates the content validity of an instrument by asking members of the target population and professionals to
rate its items \citep{terwee2018}, and reviews of AI-related instruments have applied COSMIN to studies that develop or
validate an instrument \citep{lintner2024}. This review codes how instruments are used and reported in studies that
apply them. The unit of evaluation differs, and the item-content coding of this review (C2b) is a researcher
classification of item wording, not a COSMIN rating.

\subsection{Reviews of measurement practice}
Reviews of measurement practice in user-experience research have counted item disclosure, adaptation and reasons for
choosing scales \citep{perrig2024}. In research on trust in AI, one review counted that 12 of 16 studies that measured
trust with questionnaires developed their own \citep{bach2022}, and a review of highly cited studies counted that 27\%
developed their own items and 58\% used standardized questionnaires \citep{benk2024}. No count of this review is
compared with these counts. Unlike some of these codes, a citation alone does not count as a reason for choosing an
instrument in C2d.

\subsection{Methodological studies and recalculation}
Methodological studies examine how research is designed and reported \citep{mbuagbaw2020}; their reporting has its own
guideline \citep{murad2017}, and its quality has been found wanting \citep{long2024}. Recalculating reported statistics
from their inputs is established in psychology \citep{nuijten2016} and has been used outside it, for example to
recalculate ratios and P-values from reported confidence intervals in medical abstracts \citep{georgescu2018}.

\section{Method}

\subsection{Registration and reporting}
The protocol was registered with the Generalized Systematic Review Registration form \citep{vandenakker2023} on
\tbd{date} (\tbd{identifier}), with the C2b codebook (version 2), the synthesis script and its test generator, the
screening order and the manuscript outline as attachments. Amendments are listed in Table~2 with their dates, reasons
and affected items. The report follows the checklist for meta-epidemiological methodology research
\citep{murad2017}, supplemented by the items of PRISMA 2020 that it lacks \citep{page2021}; the broader guideline for
methodological studies \citep{lawson2020} was still under development at registration.

\subsection{Source of records}
Records come from a literature database built by both authors from free sources: an exploratory bibliography of 34
records and additional candidates. This review ran no search of its own. \tbd{Searches and their dates, from the
registration attachment.}

\subsection{Eligibility and screening}
Eligibility follows the source database. Both authors screened, and the final decision was recorded with its date.
Before registration the first author had screened 71 additional candidates at title and abstract level after reading
AI-generated preliminary appraisals; the second author had not read these appraisals. After registration the remaining
candidates are screened in a fixed order, alternating between population groups, and every decision is released.

\subsection{Sets}
The exploration set is the 10 primary studies included from the exploratory bibliography and four item-coded
instruments; the rules were built from it, and it is reported separately. The main set is the primary studies among
the additional candidates for which both screeners' final decision was recorded on or before 31 December 2026, less
the calibration records. Studies included later and before main coding form a late stratum, reported outside the
primary counts. Calibration records are drawn at random from the closed pool (seed 20261017), or from the exploration
set if the pool has 24 studies or fewer. A table based on fewer than 10 coded studies reports counts without
proportions.

\subsection{Coding}
Units are the record (C1 and C3 to C7), the fatigue instrument (C2) and the item (C2b). Missing values have their own
codes, and a missing code stands alone in a cell. Table~1 lists the seven codes and their values.
\begin{itemize}
  \item C1, level of attribution: from items that mention AI to no mention, with a representative value per instrument.
  \item C2a to C2d: origin and modification, item content (types and, for reduced interest, avoidance and futility,
  targets), reported adaptation procedures, and reported measurement properties.
  \item C3: the instrument used for the AI-side variable.
  \item C4: completeness and internal consistency, including recalculation (C4f).
  \item C5: sample overlap, without deciding sameness.
  \item C6: the title against the name of the main instrument (C6a), causal wording against design (C6b), and an
  AI-fatigue term in the title against the level of attribution (C6c).
  \item C7: three conditions on the main conclusion, each its own problem type: causal wording not supported by the
  design; an AI-fatigue term while the main instrument ties fatigue to AI only through a definition, a label, an
  interview theme or not at all; and a construct term in a different family from the instrument's name.
\end{itemize}
C4f follows the logic of statcheck \citep{nuijten2016}, which recalculates p from reported test statistics and degrees
of freedom, with four differences: coders judge each test rather than scanning the whole text; values in tables are
included; p-values that the report says are adjusted for multiple testing are not counted; and percentages and indirect
effects are added. The work is done by people. An AI assistant drafts values for C4f and for the bibliographic matches
of C5 only; the drafts are stored unopened with their SHA-256 and opened only after both coders have fixed their own
values.

\subsection{Calibration, coding and reconciliation}
\tbd{Who coded, as it happened.} By default both authors code every record independently. Coding stops calibrating when
percentage agreement reaches 80\% on every registered unit; otherwise more records are drawn, up to two further rounds.
Both default coders share the view of an earlier narrative review, so their agreement is agreement between coders who
share that view. Disagreements are discussed, and those that remain are coded unresolved.

\subsection{Synthesis}
All values are computed by the registered script from the reconciled code table. The review describes; it estimates no
effect and tests no hypothesis. For each problem type, the records that can show or not show the problem are defined in
advance; among them, a problem type ``was observed in the main set'' if at least one record shows it, ``was not observed
in the main set'' if none does, and ``could not be evaluated'' if no record can show it. Each statement gives the count,
the number of evaluable records, the number not evaluable and the number unresolved. These statements report presence
or absence in this set; they are not frequencies. Sensitivity analyses use peer-reviewed records only, remove records
with unconfirmed signs of overlap, count the secondary type of boundary items, and use each coder's values in turn.

\section{Results}
\tbd{Written only after the tables have been generated from the code table, in the registered terms.}

\begin{enumerate}
  \item Records and sets (Figure~1, Table~3).
  \item Coder agreement (Table~4).
  \item RQ1 attribution (Table~5).
  \item RQ2 instruments and items (Tables~6, 7, 8 and 15).
  \item RQ3 reporting (Table~9).
  \item RQ4 sample overlap (Table~10).
  \item RQ5 and SQ2: titles and conclusions (Table~11).
  \item Problem statements (Table~12).
  \item SQ1 by centre (supplement).
  \item Sensitivity analyses (Table~13).
  \item Exploration set and late stratum (supplement).
\end{enumerate}

\begin{table}[h]
\centering
\caption{Problem statements for the main set (shell).}
\small
\begin{tabularx}{\linewidth}{XXXcccc}
\toprule
Problem type & Code and value & Statement & Count & Evaluable & Not evaluable & Unresolved \\
\midrule
--- & --- & observed / not observed / could not be evaluated & --- & --- & --- & --- \\
\bottomrule
\end{tabularx}

\smallskip
\raggedright\footnotesize One row per registered problem type; C7 has three rows, one per condition. A note gives,
for each row, the values that made records not evaluable.
\end{table}

\section{Discussion}
\tbd{Written after the results, in four parts: what the main set shows, question by question; the relation to the
narrative review, with its observations set beside the statements for the main set; recommendations for authors and
reviewers, each tied to a code and to either an existing standard or a proposal of this scheme, given whether or not
the corresponding problem was observed; and the limitations.}

\section{Conclusions}
\tbd{The scope sentence and the registered statements of the abstract, without additional claims.}

\section*{Declarations}
\textbf{Funding.} None.
\textbf{Competing interests.} No financial interests. Both authors built the source literature database \tbd{and wrote
the earlier narrative review}.
\textbf{Ethics.} Not applicable; no human participants.
\textbf{Use of AI.} An AI assistant drafted values for C4f and C5 as described in the Method \tbd{and any other use in
preparing the manuscript}.
\textbf{Author contributions.} \tbd{to be agreed}
\textbf{Data and code availability.} The registration, the code table, the codebook (CC BY 4.0) and the scripts (MIT
licence) are in the OSF project of the registration \tbd{link}.

\bibliographystyle{apalike}
\bibliography{refs}

\end{document}
